Für ein ideales Gas mit $f$ Freiheitsgraden ist $C_v = \frac{f}{2}R$ und $C_p = C_v + R$ (FoTaBe S.~171). Sauerstoff (O$_2$) ist zweiatomig, $f = \fOO$, und hat die molare Masse $M = \MOO$ (FoTaBe S.~228). Die spezifischen Wärmen sind $c = C/M$:
\begin{align}
 c_v &= \cvOF = \frac{\fOO}{2}\cdot\frac{\ncR}{\MO} = \cvO \approx \result{\cvOP{0}{2}} \\
 c_p &= \cpOF = \left(\frac{\fOO}{2}+1\right)\cdot\frac{\ncR}{\MO} = \cpO \approx \result{\cpOP{0}{2}}
\end{align}
Die Tabelle \emph{Thermische Daten von Gasen} (FoTaBe S.~191) gibt $c_p = \cpTO$. Der berechnete Wert ist nur um
\begin{align}
 \frac{c_{p,\text{Tab}}-c_p}{c_{p,\text{Tab}}} &= \dcp \approx \dcpP{0}{1}
\end{align}
kleiner: Das Modell mit fünf Freiheitsgraden beschreibt Sauerstoff bei Zimmertemperatur sehr gut.
